\section{Del salto espectral 4 a las simetrías del Monstruo: despliegue APP
\texorpdfstring{$2\to6\to54$}{2--6--54}, retículo de Leech,
\texorpdfstring{$V^\natural$}{V natural} y corredor Moonshine}
\label{sec:cierre-equivalencias-corredor-moonshine-20260825}

El corredor no es una sucesión de coincidencias numéricas. Parte del bloque
central de APP, conserva la orientación suma--producto, se ramifica en los
acoplamientos APP--Fock y APP--Witt, y vuelve a coordinarse mediante la
gradación. El orden causal vinculante es
\[\begin{aligned}
 B_c=9P_++5P_-
 &\longrightarrow 4\longrightarrow2\longrightarrow6\longrightarrow54\\
 &\longrightarrow
 \left\{
 \begin{array}{l}
  \text{grado dos de Fock y }[q]t_3=54,\\
  \text{alineamiento APP--Witt y }N_\alpha\cong\Lambda_{24}
 \end{array}
 \right.\\
 &\longrightarrow V^\natural
 \longrightarrow
 \bigl(\operatorname{Aut},\operatorname{ch},g\mapsto T_g\bigr)
 \longrightarrow\text{Moonshine}.
\end{aligned}\]
Las dos ramas intermedias poseen codominios distintos. La llave común \(54\)
es un invariante graduado transportado por el corredor; no es una aplicación
lineal desde un escalar hacia un álgebra de operadores de vértice.

\subsection{Jerarquía tipada y \(2\) realizaciones del índice \(54\)}

\begin{theorem}[Jerarquía espectral y despliegue APP]
\label{thm:cierre-vtm-jerarquia-app}
Sea \(B_c=7I+2R=9P_++5P_-\) el bloque central de APP.\@ Si
\(M_\Pi,M_\Sigma\) son las compresiones modulares definidas, entonces
\[
 9-5=4,
 \qquad (B_c)_{i\ne j}=2,
 \qquad \sum M_\Pi-\sum M_\Sigma=51-45=6,
\]
y el despliegue sobre las nueve posiciones da
\[
 9\cdot6=459-405=54.
\]
Los cuatro enteros pertenecen, respectivamente, al salto espectral, al
acoplamiento local, a la compresión modular y a la integración bidimensional.
\end{theorem}

\begin{proof}
La descomposición espectral de \(B_c\) da \(9\) y \(5\), mientras que
\(7I+2R\) fija el coeficiente no diagonal. Las dos sumas agregadas son \(51\)
y \(45\). Cada entrada comprimida representa nueve posiciones, de modo que el
despliegue multiplica la diferencia por nueve. Ninguno de los cuatro enteros
se obtiene identificando retrospectivamente sus tipos.
\end{proof}

Fijadas la recta de orientación \(o_{\Sigma\Pi}\), la paridad de hoja y la
normalización triangular, el acoplamiento orientado
\[
 \mathcal H_{\Sigma,\Pi}^{(m)}:
 M_{\mathrm{bal}}^{C_3,-}
 \longrightarrow
 \operatorname{End}\!\bigl(\mathcal F_2(\mathcal V_m)\bigr)
 \otimes o_{\Sigma\Pi}
\]
satisface
\[
 \operatorname{Tr}_{o,\mathcal F_2}
 \left(\mathcal H_{\Sigma,\Pi}^{(m)}(\nu)\Gamma_2(g_m)\right)
 =-\frac{3m\,c(\nu)}2.
\]
Para \(m=12\) y \(c(\delta)=-3\), la traza vale \(54\). Por otro lado,
\[
 t_3(q)=q^{-1}\prod_{n\ge1}(1+q^n+q^{2n})^{-12}
       =q^{-1}-12+54q-76q^2-243q^3+\cdots,
\]
de donde \([q]t_3=54\).

\begin{theorem}[Cono de compatibilidad del índice común]
\label{thm:cierre-vtm-indice-comun-54}
Con las estructuras de partida de cada rama conservadas, se cumple
\[
 \boxed{
 D_{\mathrm{APP}}
 =\operatorname{Tr}_{o,\mathcal F_2}
   \left(\mathcal H_{\Sigma,\Pi}^{(12)}(\delta)\Gamma_2(g_{12})\right)
 =[q]t_3
 =v_\alpha^2
 =\frac{196884}{5\cdot729+1}
 =54.}
\]
Esta igualdad constituye un cono de evaluaciones hacia el mismo entero. No
identifica \(W_{\mathrm{APP}}\), el espacio de Fock, el retículo vecino y
\(V^\natural\).
\end{theorem}

\begin{proof}
El censo APP da \(459-405=54\); la traza orientada da \(54\); la expansión
del producto ciclotómico da \([q]t_3=54\); y el vector radial del vecino
satisface \(v_\alpha^2=4^2\cdot2+11\cdot2=54\). Finalmente,
\(5\cdot729+1=3646\) y \(54\cdot3646=196884\). Las igualdades se comparan
después de construir cada objeto y por eso no son una selección numérica del
corredor.
\end{proof}

\paragraph{Estatuto exacto de la jerarquía \(4\to2\to6\to54\) y de la identificación reticular.}
La jerarquía \(4\to2\to6\to54\), la traza APP--Fock y la extracción
\([q]t_3\) son identidades internas exactas. La identificación del vecino con
Leech usa la cámara, el vector de vecindad y la clasificación reticular
declaradas.
%
